% =========================================================
% 6. Advisor Guidance Needs
% =========================================================
\section*{六、需要指導教授指導內容}
\addcontentsline{toc}{section}{六、需要指導教授指導內容}

本研究屬於「可部署之系統研究 + 可量測之評估設計」，需要指導教授在研究定位、方法嚴謹性、實驗設計與倫理合規等方面提供指導，以確保成果具備學術貢獻與可重現性。具體需要指導內容如下：

\subsection*{(一) 研究定位與問題定義}
\begin{itemize}
  \item 協助確認「長期指導一致性」之學術範疇與可驗證定義（研究脈絡一致性、人設一致性之界線與可量測化）。
  \item 指導研究貢獻主張的精確表述：本研究相較既有教育導師、記憶系統、人設控制方法的差異化立足點與可比較基線設計。
\end{itemize}

\subsection*{(二) 系統設計與實作決策審查}
\begin{itemize}
  \item 指導 Dependency Graph 的 schema 設計（節點/邊類型、版本化策略、衝突/修正的標記規則），以及寫入與更新邏輯是否符合「研究會議」實務。
  \item 指導多層次記憶協調策略（Working/Short-term/Long-term 的分工、淘汰策略、檢索排序與噪聲控制），避免記憶錯誤累積導致的系統性偏誤。
  \item 指導 Persona Vector 的抽取流程與干預層選擇（中後層範圍、監控分數聚合方式），並確認干預強度 $\alpha$ 的安全範圍與失效模式（過度干預導致僵化/離題）。
\end{itemize}

\subsection*{(三) 實驗設計、統計檢定與可重現性}
\begin{itemize}
  \item 協助審查壓力測試腳本的有效性：矛盾挑戰是否真能測出跨會議一致性差異；漂移量化指標是否可區分「風格波動」與「本質人設漂移」。
  \item 指導真人盲測的實驗設計：受試者招募條件、樣本數估計、within-subject 隨機化、干擾因子控制、統計方法（例如配對檢定/重複量數 ANOVA/混合效應模型擇一）。
  \item 指導評分規準與量表使用：MCA-21 的題項詮釋、計分與信度檢查（例如一致性、評分者間一致性），以及如何與系統自動化指標對齊。
\end{itemize}

\subsection*{(四) 研究倫理、資料治理與風險控管}
\begin{itemize}
  \item 指導 IRB/倫理審查需求：真人實驗的同意程序、匿名化、資料保存期限與可撤回機制。
  \item 指導學術情境中的安全性邊界：避免系統產生不當建議（例如過度自信、缺乏證據的結論），並設計必要的警示與不確定性表達機制。
  \item 協助規劃地端部署/隱私需求的論述：在校園或研究室環境中如何降低敏感研究資料外流風險。
\end{itemize}

\subsection*{(五) 進度控管與論文撰寫}
\begin{itemize}
  \item 定期檢視里程碑產出（系統原型、圖式記憶更新規則、人設監控儀表板、壓力測試報告、真人實驗結果）。
  \item 指導成果撰寫：研究方法敘述的可重現性、消融實驗（ablation）與負結果呈現方式、威脅效度（threats to validity）與限制的嚴謹論證。
\end{itemize}
